\documentclass[12pt,a4paper]{report}

\usepackage[T1]{fontenc}
\usepackage[utf8]{inputenc}
\usepackage{lmodern}
\usepackage{geometry}
\usepackage{fancyhdr}
\usepackage{xcolor}
\usepackage{hyperref}
\usepackage{booktabs}
\usepackage{array}
\usepackage{tabularx}
\usepackage{longtable}
\usepackage{setspace}
\usepackage{verbatim}

\geometry{margin=2.5cm}
\setlength{\headheight}{15pt}
\setstretch{1.15}

\renewcommand{\contentsname}{Table des matieres}
\renewcommand{\chaptername}{Chapitre}

\definecolor{titleblue}{HTML}{1F4E79}
\definecolor{lightblue}{HTML}{EAF2F8}
\definecolor{codegray}{HTML}{F6F8FA}
\definecolor{darkgray}{HTML}{444444}

\hypersetup{
    colorlinks=true,
    linkcolor=titleblue,
    urlcolor=titleblue,
    pdftitle={Rapport de projet - Systeme de gestion de circulation},
    pdfauthor={Projet universitaire}
}

\pagestyle{fancy}
\fancyhf{}
\fancyhead[L]{Systeme de gestion de circulation}
\fancyhead[R]{Rapport de projet}
\fancyfoot[C]{\thepage}

\begin{document}

\begin{titlepage}
    \centering
    \vspace*{2cm}

    {\Huge\bfseries\color{titleblue} Rapport de projet\par}
    \vspace{0.6cm}
    {\LARGE Systeme distribue de supervision d'un carrefour intelligent\par}

    \vspace{1.5cm}
    \rule{0.9\textwidth}{1.2pt}\par
    \vspace{0.4cm}
    {\Large Simulation de capteurs, services SOAP, broker Kafka, persistance MySQL et API JAX-RS\par}
    \vspace{0.4cm}
    \rule{0.9\textwidth}{1.2pt}\par

    \vspace{2cm}

    \begin{tabular}{>{\bfseries}p{4cm}p{8cm}}
        Projet : & Gestion de la circulation sur un carrefour a quatre voies \\
        Technologies : & Java, JAX-WS, Axis2, Kafka, MySQL, JAX-RS, Tomcat \\
        Modules : & \texttt{client}, \texttt{services}, \texttt{centraleservice} \\
        Type : & Projet universitaire \\
        Date : & \today \\
    \end{tabular}

    \vfill

    \begin{center}
        \fcolorbox{titleblue}{lightblue}{
            \parbox{0.82\textwidth}{
                \centering
                Ce document presente l'architecture, le fonctionnement et les choix techniques d'une plateforme de simulation de trafic routier, avec deux services principaux : \texttt{ServiceFlux} et \texttt{ServiceFeux}.
            }
        }
    \end{center}

    \vfill
    {\large Annee universitaire 2025--2026\par}
\end{titlepage}

\tableofcontents
\clearpage

\chapter*{Resume}
\addcontentsline{toc}{chapter}{Resume}
Ce projet realise une chaine complete de collecte, de transport, de stockage et d'exploitation de donnees de circulation dans le cadre d'un carrefour intelligent simule. La solution s'appuie sur trois sous-systemes principaux. Le module \texttt{client} simule quatre voies de circulation et l'etat des feux. Le module \texttt{services} expose deux services SOAP responsables de la reception des mesures et de la transmission de l'etat logique du carrefour. Enfin, le module \texttt{centraleservice} consomme les messages de flux depuis Kafka, les persiste dans MySQL et met a disposition une API JAX-RS consultee par une interface web.

L'objectif pedagogique du projet est double. D'une part, il s'agit de mettre en oeuvre plusieurs technologies de middleware dans une meme application coherente. D'autre part, il s'agit de proposer une architecture suffisamment modulaire pour separer l'acquisition des donnees, la gestion des etats du carrefour et la consultation par un systeme central. Cette separation facilite la maintenance, la demonstration et l'evolution du projet.

\chapter{Introduction generale}
La gestion de la circulation urbaine constitue un cas d'etude interessant pour les architectures distribuees, car elle implique des sources de donnees multiples, des evenements quasi temps reel, des traitements differencies et des besoins de supervision centralisee. Dans ce projet, le carrefour a quatre voies sert de support a une simulation de trafic dans laquelle chaque route produit periodiquement des informations de flux. En parallele, les feux du carrefour possedent un etat logique partage qui influence le comportement des vehicules simules.

Le projet a ete structure en trois modules afin de bien repartir les responsabilites. Le \texttt{client} represente la partie terrain. Il genere les mesures de trafic et construit l'etat courant du carrefour. Le module \texttt{services} joue le role d'intermediaire SOAP. Il recoit les informations des capteurs et les transmet vers les couches superieures. Enfin, le module \texttt{centraleservice} joue le role de service central : il conserve l'historique de flux, signale les congestions, expose des endpoints REST et permet a une application front-end de consulter ou de modifier l'etat de l'intersection.

Le choix de cette architecture permet de combiner plusieurs approches complementaires. SOAP est utilise entre les simulateurs et les services metiers, car il fournit un contrat explicite a travers le WSDL et illustre bien les notions de stub, de port et d'interoperabilite entre frameworks. Kafka est employe comme broker pour decoupler la collecte du flux et sa persistance. Enfin, JAX-RS offre une interface plus legere et plus moderne pour la consultation depuis le front-end. Cette coexistence de styles de communication correspond bien aux objectifs d'un projet universitaire, car elle met en valeur plusieurs paradigmes dans un meme systeme.

D'un point de vue fonctionnel, la plateforme doit repondre a plusieurs besoins. Elle doit recuperer les mesures de circulation sur chaque route, stocker ces mesures en base de donnees, detecter les congestions selon un seuil donne, fournir l'etat des feux, autoriser le changement de duree des cycles et permettre le for\c{c}age de certaines orientations du carrefour. Le tout doit rester suffisamment simple pour etre deploye, teste et explique dans un contexte d'evaluation.

\chapter{Architecture generale du systeme}
L'architecture du projet s'articule autour de trois briques principales reliees par des protocoles differents. La figure logique ci-dessous peut etre decrite textuellement de la maniere suivante :

\begin{enumerate}
    \item Le module \texttt{client} simule quatre capteurs de flux et un emetteur d'etat des feux.
    \item Les capteurs de flux appellent \texttt{ServiceFlux} en SOAP selon une approche JAX-WS.
    \item L'etat des feux est envoye a \texttt{ServiceFeux} via un stub Axis2.
    \item \texttt{ServiceFlux} publie les mesures dans Kafka.
    \item Le module \texttt{centraleservice} lance un consommateur Kafka qui lit les mesures et les stocke dans MySQL.
    \item En parallele, \texttt{ServiceFeux} transmet les mises a jour de l'etat des feux a la centrale et accepte egalement des commandes en sens inverse.
    \item L'application web consomme finalement l'API JAX-RS de la centrale pour afficher l'etat du systeme et envoyer des actions de pilotage.
\end{enumerate}

\section*{Organisation des modules}
\begin{longtable}{p{3cm}p{10.5cm}}
\toprule
\textbf{Module} & \textbf{Role principal} \\
\midrule
\endhead
\texttt{client} & Simulation des routes, calcul du flux courant, evolution du temps restant des feux, envoi des donnees de flux et de l'etat du carrefour. \\
\texttt{services} & Exposition des services SOAP \texttt{ServiceFlux} et \texttt{ServiceFeux}, publication vers Kafka et relais de certaines informations vers la centrale. \\
\texttt{centraleservice} & Consommation Kafka, persistance MySQL, gestion des alertes, exposition des endpoints JAX-RS et commande du carrefour depuis la couche centrale. \\
\bottomrule
\end{longtable}

\section*{Choix techniques}
Le projet ne repose pas sur une seule technologie, mais sur une combinaison volontairement representative d'un systeme distribue :

\begin{itemize}
    \item \textbf{JAX-WS} pour la communication entre les capteurs de flux et \texttt{ServiceFlux}.
    \item \textbf{Axis2} pour le client de communication avec \texttt{ServiceFeux}, ce qui montre qu'un client Axis2 peut interoperer avec un service SOAP expose cote serveur en JAX-WS.
    \item \textbf{Kafka} pour decoupler l'acquisition des mesures de trafic de leur stockage.
    \item \textbf{MySQL} pour conserver l'historique des flux et permettre les consultations ulterieures.
    \item \textbf{JAX-RS} pour fournir une API facile a interroger depuis un front-end web.
    \item \textbf{Tomcat} pour le deploiement du service central.
\end{itemize}

L'utilisation de ces choix est pedagogiquement pertinente. Elle permet de montrer qu'un systeme peut melanger plusieurs styles de communication sans perdre en coherence, a condition de bien distinguer les responsabilites. Le flux monte du terrain vers le broker et la base de donnees, tandis que les commandes de la centrale vers les feux suivent un chemin de retour maitrise.

\section*{Scenario nominal d'execution}
Lorsqu'on lance le projet dans des conditions normales, le deroulement est le suivant. Les threads de simulation sont demarres dans le \texttt{client}. Chaque capteur de flux modifie localement la valeur de trafic de sa route puis envoie une mesure toutes les cinq secondes a \texttt{ServiceFlux}. Le service construit un objet JSON contenant le nom de la route, la valeur du flux et l'horodatage, puis publie ce message dans le topic Kafka dedie. En arriere-plan, le \texttt{centraleservice} consomme les messages du topic, les transforme en objets metier et les persiste dans MySQL. L'API REST peut alors retourner l'historique, la derniere mesure par route ou encore les alertes de congestion.

En parallele, le simulateur des feux envoie regulierement l'etat logique du carrefour a \texttt{ServiceFeux}. Celui-ci publie aussi l'information dans Kafka et la relaie a la centrale. Quand l'utilisateur modifie la duree ou force une orientation via l'API REST, la centrale appelle \texttt{ServiceFeux} pour transmettre une commande de retour. Cette commande est ensuite lue par les capteurs pour ajuster leur etat partage. L'intersection reste donc coherente, car le changement applique a une voie impacte la logique globale du carrefour.

\chapter{Le service \texttt{ServiceFlux}}
Le service \texttt{ServiceFlux} constitue la porte d'entree des mesures de trafic dans l'architecture. Son role est simple en apparence, mais il est essentiel car il relie la simulation terrain au pipeline de traitement central. Le capteur de flux de chaque route simule l'evolution du nombre de vehicules en fonction de l'etat des feux et publie regulierement cette valeur vers le service. Dans l'implementation retenue, le client utilise un stub JAX-WS genere a partir du WSDL local, ce qui garantit une integration claire et relativement stable entre le simulateur et le service expose.

Sur le plan technique, \texttt{ServiceFlux} est un endpoint SOAP publie en local a l'adresse \texttt{http://localhost:8080/ServiceFlux}. Il implemente l'interface \texttt{IServiceFlux} et expose une operation principale, \texttt{sendFlux(int flux, String name)}. Cette operation recoit la valeur du trafic et l'identifiant de la route, puis enrichit ces donnees avec un horodatage genere cote service. Ce choix evite de dependre entierement du temps cote client et permet d'associer a chaque message une date de reference unifiee au moment de son passage dans la couche metier.

Une fois les donnees recues, le service ne tente pas de les stocker directement en base. A la place, il adopte une logique de publication asynchrone dans Kafka. La mesure est serialisee en JSON avec les champs \texttt{flux}, \texttt{name} et \texttt{timestamp}, puis envoyee au broker via l'utilitaire \texttt{KafkaProducerUtil}. Cette etape est tres importante dans l'architecture. Elle evite que le service SOAP soit directement couple a la base de donnees ou au service central. Le service reste donc leger, focalise sur la reception des mesures et la mise a disposition des evenements.

Le recours a Kafka apporte plusieurs avantages. D'abord, il introduit un decouplage temporel : le capteur n'a pas besoin d'attendre que la centrale traite immediatement la donnee. Ensuite, il permet d'imaginer une evolution du projet dans laquelle plusieurs consommateurs pourraient lire les memes mesures, par exemple pour de l'archivage, de l'analyse statistique ou de la visualisation en temps reel. Enfin, il reduit l'impact d'une indisponibilite temporaire de la base de donnees sur le chemin d'acquisition. Tant que le broker est disponible, la mesure peut etre publiee et consommee ulterieurement.

Le service s'insere donc dans une chaine de responsabilites bien precise. Le capteur calcule localement un flux simule. \texttt{ServiceFlux} le recoit, le valide implicitement par son typage et le transforme en evenement. Le consommateur Kafka du module central lit ensuite cet evenement, le deserialize en \texttt{ReceptionFlux} et invoque \texttt{GestionFlux.insertData(...)} afin de persister la mesure dans MySQL. La logique de stockage n'est donc pas dupliquee dans plusieurs couches. Elle est centralisee dans le service central, ce qui simplifie la maintenance.

Du point de vue fonctionnel, \texttt{ServiceFlux} participe indirectement a deux besoins majeurs : l'historisation et la detection des alertes. L'historique est rendu accessible par les endpoints REST \texttt{/Flux}, \texttt{/Flux/latest} et \texttt{/Flux/route/\{name\}}. Les alertes sont construites lorsque la mesure depasse le seuil de congestion defini dans \texttt{GestionFlux}. Meme si le service ne declenche pas directement l'alerte, il est le point de depart de cette logique, car toute la chaine est alimentee par les mesures qu'il publie.

Le choix de JAX-WS pour ce service est coherent avec la nature du projet. Il permet de generer facilement le client a partir du WSDL et de montrer le fonctionnement classique des services web SOAP. Dans ce contexte universitaire, \texttt{ServiceFlux} joue donc un double role. Il est a la fois un composant utile du point de vue metier et un support concret pour illustrer la notion de contrat de service, de stub genere et d'integration distribuee. Son implementation reste volontairement simple, mais elle remplit pleinement sa fonction dans le systeme global.

\section*{Extrait representatif}
\begin{verbatim}
@WebService(
    targetNamespace = "http://Services.hub.service/",
    portName = "ServiceFluxPort",
    serviceName = "ServiceFlux",
    endpointInterface = "service.hub.Interfaces.IServiceFlux"
)
public class ServiceFlux implements IServiceFlux {
    @Override
    public void sendFlux(int flux, String name) {
        String timestamp = LocalDateTime.now()
            .format(DateTimeFormatter.ofPattern("yyyy-MM-dd HH:mm:ss"));

        String json = String.format(
            "{\"flux\":%d,\"name\":\"%s\",\"timestamp\":\"%s\"}",
            flux, name, timestamp
        );
        KafkaProducerUtil.publish(name, json);
    }
}
\end{verbatim}

\chapter{Le service \texttt{ServiceFeux}}
Le service \texttt{ServiceFeux} remplit une fonction plus riche que \texttt{ServiceFlux}, car il ne se contente pas de remonter de l'information : il participe egalement au pilotage du carrefour. Son premier role consiste a recevoir l'etat courant des feux calcule par la simulation. Toutes les secondes, le composant \texttt{CapteurFeux} construit un tableau de quatre objets representant les voies du carrefour, avec le nom visible de la route, le segment logique auquel elle appartient et le temps restant avant changement. Ce tableau est envoye a \texttt{ServiceFeux} via un stub Axis2.

Ce choix est interessant du point de vue de l'interoperabilite. Le client cote simulation utilise Axis2, tandis que le service lui-meme est publie cote serveur avec JAX-WS. Cette combinaison montre qu'un meme contrat SOAP peut etre consomme par des outils differents tant que le WSDL reste compatible. Pour un projet universitaire, cette heterogeneite est plutot un avantage, car elle permet de demontrer la comprehension des mecanismes de generation de stubs, de ports et d'invocation de services web.

La methode \texttt{setFeux(Feux[] list)} est le point d'entree principal pour la remontee d'etat. Elle transforme le tableau de feux en JSON, publie cette representation dans Kafka avec la cle \texttt{feux} et transmet la meme information a la centrale via un appel HTTP vers l'endpoint \texttt{/Feux/maj}. La centrale met alors a jour sa vision du carrefour. Cette vision inclut non seulement la couleur logique de chaque orientation, mais aussi le temps restant, les identifiants stables de route et les noms visibles affiches par le front-end. La centralisation de cet etat permet au front-end de consulter l'etat courant sans avoir a parler directement aux capteurs ou aux services SOAP.

Le deuxieme role de \texttt{ServiceFeux} concerne le chemin inverse, c'est-a-dire la commande du carrefour depuis la centrale. Pour cela, deux operations supplementaires ont ete introduites : \texttt{setFeuxTemp(String name, int feux)} et \texttt{getFeuxTemp(String name)}. Le principe retenu est le suivant. Lorsqu'un utilisateur modifie la duree des feux ou force une orientation a travers l'API JAX-RS, le module central appelle \texttt{setFeuxTemp}. Une commande est alors enregistree dans le service avec une duree et une orientation logique. Les capteurs de flux, qui executent la simulation au fil du temps, consultent regulierement \texttt{getFeuxTemp(name)}. Si une nouvelle commande existe pour leur route, ils l'appliquent a l'etat partage de l'intersection.

Cette approche presente un avantage important : elle conserve une logique de carrefour globale. Dans l'implementation du projet, une voie n'est pas pilotee isolement. Une modification sur \texttt{nord} ou \texttt{sud} agit en realite sur le segment nord-sud, tandis que \texttt{est} et \texttt{ouest} constituent le segment oppose. Ainsi, lorsqu'une route est forcee au vert, sa route opposee le devient aussi, et l'autre segment passe au rouge. Cette contrainte correspond mieux a la realite d'un carrefour a deux axes qu'a une gestion totalement independante de chaque feu.

Le service \texttt{ServiceFeux} sert donc de point de synchronisation entre trois mondes : la simulation locale, la centrale REST et les mecanismes d'integration. Il met a jour la centrale en continu, mais il permet aussi a la centrale de reconfigurer dynamiquement le carrefour. Cette bidirectionnalite est l'un des aspects les plus interessants du projet, car elle montre qu'un service web peut etre utilise a la fois comme canal de remontee d'etat et comme canal de commande.

Enfin, ce service rend possibles plusieurs fonctionnalites visibles par l'utilisateur final. Le front-end peut consulter \texttt{/Feux}, \texttt{/Feux/etat} et \texttt{/Feux/config} pour suivre l'etat du carrefour. Il peut egalement appeler \texttt{/Feux/config}, \texttt{/Feux/force/\{name\}} ou \texttt{/Feux/nom/\{routeId\}} pour agir sur la duree, sur la direction active ou sur le nom affichable d'une route. Sans \texttt{ServiceFeux}, ces operations resteraient purement locales a la centrale. Grace a lui, elles peuvent se propager jusque dans la simulation.

\section*{Extrait representatif}
\begin{verbatim}
@WebMethod(operationName = "setFeuxTemp")
public synchronized boolean setFeuxTemp(
    @WebParam(name = "name") String name,
    @WebParam(name = "feux") int feux
) {
    int nextDuration = Math.max(1, Math.abs(feux));
    boolean routeShouldBeGreen = feux >= 0;
    boolean routeOnNorthSouth = isNorthSouth(name);

    feuxDuration = nextDuration;
    northSouthGreen = routeShouldBeGreen
        ? routeOnNorthSouth
        : !routeOnNorthSouth;
    commandVersion++;
    return true;
}
\end{verbatim}

\chapter{Le service central JAX-RS et la persistance}
Le module \texttt{centraleservice} constitue la couche de supervision du projet. Il est deployee sur Tomcat et expose une API JAX-RS qui sert d'interface de consultation et de commande pour un front-end web. Sa premiere responsabilite est de persister les mesures de flux dans MySQL. Cette tache est assuree par un consommateur Kafka lance au demarrage de l'application via un \texttt{ServletContextListener}. Le thread de consommation reste actif tant que l'application est deployee et s'arrete proprement lors de l'arret du contexte.

Le consommateur lit les messages du topic configure, ignore ceux dont la cle correspond a \texttt{feux} et deserialize les autres sous forme d'objets \texttt{ReceptionFlux}. Chaque mesure est ensuite stockee via \texttt{GestionFlux.insertData(...)}. Cette methode ajoute aussi une alerte en memoire lorsque le flux depasse le seuil de congestion. La centrale joue ainsi un double role : base d'historisation et detecteur simple d'evenements critiques.

Au-dessus de cette couche de persistance, l'API JAX-RS fournit plusieurs familles d'endpoints. Les endpoints \texttt{/Flux} retournent l'historique general ou par route. L'endpoint \texttt{/Alert} retourne les alertes memorisees depuis la derniere consultation. Les endpoints \texttt{/Feux} et \texttt{/Feux/etat} donnent l'etat courant detaille des feux. Enfin, des endpoints d'ecriture permettent de modifier la configuration, de forcer une orientation ou de renommer une route. L'ensemble constitue une interface suffisamment complete pour etre integree a une application web de demonstration.

Le choix de JAX-RS pour cette couche est judicieux, car il permet d'offrir des URLs simples, des echanges JSON directs et une integration naturelle avec le front-end. Alors que les communications SOAP sont utiles pour la partie pedagogique et l'interoperabilite des services, la couche REST repond mieux au besoin de supervision interactive. Cette combinaison illustre bien la difference entre une interface de service metier formalisee par WSDL et une interface de consultation/pilotage plus legere.

\chapter{Conclusion}
Le projet realise une architecture distribuee complete autour d'un carrefour intelligent simule. Il combine la generation de donnees par des capteurs, l'utilisation de services web SOAP, la mediation par Kafka, la persistance dans MySQL et l'exposition d'une API REST exploitable par une interface web. Cette combinaison en fait un bon projet de synthese pour mettre en pratique plusieurs notions vues en genie logiciel et en systemes distribues.

Les deux services principaux, \texttt{ServiceFlux} et \texttt{ServiceFeux}, remplissent des roles complementaires. Le premier transporte les mesures de trafic vers le pipeline de stockage. Le second represente l'etat dynamique du carrefour et permet egalement son pilotage depuis la centrale. Ensemble, ils assurent la liaison entre la simulation et la supervision.

Bien que certaines ameliorations restent envisageables, comme l'externalisation plus poussee de la configuration ou la reduction de certaines dependances locales, l'etat actuel du projet est tout a fait coherent pour une livraison universitaire. L'architecture est comprensible, les composants sont identifies clairement, et les fonctionnalites principales attendues sont couvertes.

\appendix
\chapter{Exemples d'API disponibles dans la centrale}
\begin{itemize}
    \item \texttt{GET /centrale/api/Flux}
    \item \texttt{GET /centrale/api/Flux/latest}
    \item \texttt{GET /centrale/api/Flux/route/nord}
    \item \texttt{GET /centrale/api/Alert}
    \item \texttt{GET /centrale/api/Feux/etat}
    \item \texttt{GET /centrale/api/Feux/config}
    \item \texttt{POST /centrale/api/Feux/config}
    \item \texttt{POST /centrale/api/Feux/force/nord}
    \item \texttt{POST /centrale/api/Feux/nom/nord}
\end{itemize}

\end{document}
